% =========================================================
% المعوقات والأولويات — ARABIC (Section II)
% =========================================================

\newgeometry{right=2cm, left=2cm, top=3cm, bottom=2.5cm}

\noindent
\begin{minipage}[t]{0.80\textwidth}
    \vspace{0pt}
    {\fontsize{18}{22}\selectfont \textbf{\color{ecestitle} ثانياً: المعوقات التي واجهت مجتمع الأعمال خلال الربع محل الدراسة، وأولويات تحسين بيئة الأعمال من منظور عينة الشركات}}
\end{minipage}%
\hfill
\begin{minipage}[t]{0.18\textwidth}
    \vspace{0pt}
    \raggedleft
    
\includegraphics[width=\linewidth]{ph.png}
\end{minipage}

\vspace{1.5em}
\label{sec:ar_constraints}

\noindent
{\fontsize{11}{14}\selectfont \textbf{\color{ecestitle} \textit{2-1 المعوقات التي واجهت مجتمع الأعمال خلال الربع محل الدراسة}}}

\vspace{0.8em}

\noindent
{\fontsize{10}{13}\selectfont \textbf{\ECESContent{constraints_intro_ar}{[TODO: أدخل أبرز المعوقات التي تصدرت قائمة التحديات لجميع الشركات خلال ربع \CurrentQuarterText.]}}}

\vspace{0.8em}

\noindent
{\fontsize{10}{13}\selectfont
\ECESContent{constraints_main_text_p1_ar}{[TODO: أدخل جملة تشير إلى الشكل 2-1 الذي يعرض قائمة المعوقات المرتبة حسب شدتها لربع \CurrentQuarterText.]}

\vspace{0.8em}

\ECESContent{constraints_main_text_p2_ar}{[TODO: أدخل شرحاً تفصيلياً لأبرز 5 معوقات وتأثيرها على مجتمع الأعمال خلال \CurrentQuarterText.]}
}

\clearpage

%% TODO: Upload ch13.png (Arabic) — قائمة المعوقات المرتبة (شكل 2-1) via Chart Manager
\begin{figure}[h!]
    \centering
    \includegraphics[width=\linewidth]{ch13.png}
\end{figure}

\noindent
{\fontsize{8}{10}\selectfont
\ECESContent{constraints_source_pg17_ar}{المصدر: نتائج المسح.\\
\textit{ملاحظة: بعض المؤشرات مستحدثة اعتباراً من الربع محل الدراسة. الأرقام تمثل نسبة الشركات التي أشارت إلى كل معوق، ويمكن لكل شركة اختيار أكثر من معوق.}}
}

\vspace{1.5em}

\noindent
{\fontsize{11}{14}\selectfont \textbf{\color{ecestitle} \textit{2-1-1 المعوقات وفقاً لحجم الشركات}}}
\vspace{0.8em}

\noindent
{\fontsize{10}{13}\selectfont
\ECESContent{constraints_by_size_p1_ar}{[TODO: أدخل جملة تمهيدية تشير إلى أن الشركات الصغيرة والمتوسطة تواجه عموماً معوقات أكثر من الشركات الكبيرة خلال \CurrentQuarterText.]}

\vspace{0.8em}

\ECESContent{constraints_by_size_p2_ar}{[TODO: أدخل قائمة المعوقات المرتبة للشركات الكبيرة خلال \CurrentQuarterText.]}

\vspace{0.8em}

\ECESContent{constraints_by_size_p3_ar}{[TODO: أدخل قائمة المعوقات المرتبة للشركات الصغيرة والمتوسطة خلال \CurrentQuarterText.]}

\vspace{0.8em}

\ECESContent{constraints_by_size_p4_ar}{[TODO: أدخل جملة إحالة إلى الشكل 2-2.]}
}

\clearpage

%% TODO: Upload ch14.png (Arabic) — المعوقات وفقاً لحجم الشركة (شكل 2-2) via Chart Manager
\begin{figure}[h!]
    \centering
    \includegraphics[width=\linewidth]{ch14.png}
\end{figure}

\noindent
{\fontsize{8}{10}\selectfont
\ECESContent{constraints_source_pg18_ar}{المصدر: نتائج المسح.}
}

\vspace{1.5em}

\noindent
{\fontsize{11}{14}\selectfont \textbf{\color{ecestitle} \textit{2-1-2 المعوقات وفقاً للقطاعات الاقتصادية}}}
\vspace{0.8em}

\noindent
{\fontsize{10}{13}\selectfont
\ECESContent{constraints_by_sector_ar}{[TODO: أدخل تحليل تباين المعوقات عبر القطاعات الاقتصادية خلال \CurrentQuarterText — أبرز معوق لكل قطاع (شكل 2-3).]}
}

\clearpage

%% TODO: Upload ch15.png (Arabic) — المعوقات وفقاً للقطاعات (شكل 2-3) via Chart Manager
\begin{figure}[h!]
    \centering
    \includegraphics[width=\linewidth]{ch15.png}
\end{figure}

\noindent
{\fontsize{8}{10}\selectfont
\ECESContent{constraints_source_pg19_ar}{\textit{المصدر:} نتائج المسح.}
}

\clearpage

% =========================================================
% أولويات تحسين بيئة الأعمال — ARABIC
% =========================================================

\newgeometry{right=2cm, left=2cm, top=2.5cm, bottom=2.5cm}

\noindent
\begin{minipage}[t]{0.75\textwidth}
    \vspace{0pt}
    {\fontsize{20}{24}\selectfont \textbf{\color{ecestitle} 2-2 أولويات تحسين بيئة الأعمال في مصر (وفقاً لآراء عينة الشركات)}\par}
\end{minipage}%
\hfill
\begin{minipage}[t]{0.22\textwidth}
    \vspace{0pt}
    \raggedleft
    
\includegraphics[width=\linewidth, keepaspectratio]{ph2.png}
\end{minipage}

\vspace{2.5em}

\noindent
{\fontsize{10}{13}\selectfont \textbf{\ECESContent{priorities_headline_ar}{[TODO: أدخل جملة تمهيدية مميزة تلخص أبرز 3-4 أولويات لتحسين بيئة الأعمال.]}}}

\vspace{1em}

\noindent
{\fontsize{10}{13}\selectfont
\ECESContent{priorities_all_firms_ar}{[TODO: أدخل الأولويات المرتبة لجميع الشركات خلال \CurrentQuarterText.]}

\vspace{0.8em}

\ECESContent{priorities_large_firms_ar}{[TODO: أدخل الاختلافات في ترتيب الأولويات بين الشركات الكبيرة والصغيرة والمتوسطة.]}
}

\clearpage

%% TODO: Upload ch16-ch19.png (Arabic) — أشكال الأولويات via Chart Manager
\begin{figure}[h!]
    \centering
    \includegraphics[width=\linewidth]{ch16.png}
\end{figure}
\noindent
{\fontsize{8}{10}\selectfont \ECESContent{priorities_source_pg20_ar}{المصدر: نتائج المسح.}}
\clearpage

\begin{figure}[h!]
    \centering
    \includegraphics[width=\linewidth]{ch17.png}
\end{figure}
\noindent
{\fontsize{8}{10}\selectfont \ECESContent{priorities_source_pg21_ar}{المصدر: نتائج المسح.}}
\clearpage

\begin{figure}[h!]
    \centering
    \includegraphics[width=\linewidth]{ch18.png}
\end{figure}
\noindent
{\fontsize{8}{10}\selectfont \ECESContent{priorities_source_pg22_ar}{المصدر: نتائج المسح.}}
\clearpage

\begin{figure}[h!]
    \centering
    \includegraphics[width=\linewidth]{ch19.png}
\end{figure}
\noindent
{\fontsize{8}{10}\selectfont \ECESContent{priorities_source_pg23_ar}{المصدر: نتائج المسح.}}
\clearpage
